\subsection{LCD表示の処理}
\label{sec:lcd-impl}

表示用の子機は，16文字2行のLCDキャラクタディスプレイ1枚に，左右のカエル2体と中央の歌詞を同時に表示する．表示例は図\ref{fig:lcd}に示したとおりである．カエルはHD44780のカスタム文字機能で作成した5$\times$8ドットのパターンを3列2行に並べて構成し，上段の3文字は固定の絵柄，下段の3文字を口の開いた絵柄と閉じた絵柄とで差し替えることで口パクを表現する．左のカエルを列0〜2に，右のカエルを列13〜15に配置し，その間の列3〜12の2行を歌詞表示エリアとした．HD44780が同時に登録できるカスタム文字は8種類までであるため，左右のカエルで同じ6種類の文字を共有し，2体の口が常に同じ形で動く構成としている．

表示の進行は親機の赤外線tickに同期させる．受信したマスクのうち担当ビットが立ったtickでのみ歌詞位置\texttt{localPos}を1つ進め，対応する語を表示すると同時に口の開閉を反転させる．歌詞は楽譜1周分の32\,tickと1対1に対応する文字列テーブルとして保持している．1語を書き込むたびに書き込み列を語長だけ進め，残り幅に収まらなくなったら下段へ改行し，下段も埋まったら歌詞エリアだけを空白で消去して先頭から書き直す．画面全体を消去せずエリアを限定して消すのは，カエルの絵柄を保持したまま歌詞だけを更新するためである．口パクの更新では，CGRAMを書き換えると表示中の文字の形も即座に変わるHD44780の性質を利用し，カーソルを左右のカエルの下段へ移動して6文字を書き直すだけで2体の口を同時に動かしている．主要な関数を表\ref{tab:lcd_func}に，処理の流れを図\ref{fig:flow_lcd}に示す．

表示子機が受信するフレームは演奏子機と同じものであり，\ref{sec:proto-compare}節で述べたSony SIRC 12ビットへの移行にあわせて受信処理も更新している．親機は参加中の子機が1台でもあれば毎tickフレームを送信するので，表示子機は演奏子機と同じ時間基準で歌詞と口パクを進められる．ただし表示子機も演奏子機と同様に受信回数で位置を数えるため，フレームを取りこぼすと歌詞の位置が1つ後方へずれたまま戻らない．この性質と対策は\ref{sec:disc-mask}節で演奏子機とあわせて論じる．

\begin{table}[tb]
\caption{LCD表示の主要関数}
\label{tab:lcd_func}
\begin{center}
\small
\begin{tabular}{llp{18zw}}
\hline
関数名 & 役割 & 概要 \\
\hline
\texttt{setup} & 初期化 & シリアル通信とLCDを初期化してカエルのカスタム文字6種を登録し，口を閉じた状態で描画したうえで赤外線受信を開始する \\
\texttt{loop} & 受信と進行 & 赤外線フレームを受信し，担当ビットが立っていれば歌詞位置を1つ進めて歌詞表示と口パクを実行する \\
\texttt{showLyric} & 歌詞表示 & 指定位置の語を歌詞エリアへ書き込み，行幅に収まらない場合は改行またはエリアの消去を行う \\
\texttt{clearLyricArea} & 歌詞エリアの消去 & 列3〜12の2行だけを空白で埋め，書き込み位置を先頭へ戻す \\
\texttt{updateFrogs} & 口パクの更新 & 開閉に応じた口のパターンをCGRAMへ登録し直し，左右のカエル6文字を書き直す \\
\hline
\end{tabular}
\end{center}
\end{table}

\begin{figure}[tb]
\begin{center}
\begin{tikzpicture}[
  term/.style={draw, rounded corners=8pt, minimum width=9em, minimum height=2.2em, align=center, font=\footnotesize},
  proc/.style={draw, minimum width=12em, minimum height=2.2em, align=center, font=\footnotesize},
  dec/.style={draw, diamond, aspect=2.2, align=center, font=\footnotesize, inner sep=1.5pt},
  lbl/.style={font=\footnotesize},
  >=Stealth, node distance=0.55cm]
  \node[term] (start) {起動};
  \node[proc, below=of start] (init) {LCDを初期化してカエルのカスタム文字を登録し，\\口を閉じた状態で描画して赤外線受信を開始};
  \node[term, below=of init] (wait) {赤外線フレームの受信を待つ};
  \node[dec, below=of wait] (dec1) {フレームを\\受信したか};
  \node[dec, below=of dec1] (dec2) {maskの担当\\ビットが1か};
  \node[proc, below=of dec2] (pos) {歌詞位置\texttt{localPos}を1つ進める};
  \node[dec, below=of pos] (dec3) {\texttt{localPos}が\\32\,tickに達したか};
  \node[proc, below=of dec3] (lyric) {\texttt{showLyric}で歌詞エリアに次の語を書き込み\\（行幅を超える場合は改行）};
  \node[proc, below=of lyric] (flip) {口の開閉フラグ\texttt{frogOpen}を反転し，\\\texttt{updateFrogs}で左右のカエルの口を書き換える};
  \node[proc, right=0.9cm of dec3] (reset) {\texttt{localPos}を0に戻し\\歌詞エリアを消去};
  \coordinate (lane) at (6.9, 0);
  \draw[->] (start) -- (init);
  \draw[->] (init) -- (wait);
  \draw[->] (wait) -- (dec1);
  \draw[->] (dec1) -- node[lbl, right] {Yes} (dec2);
  \draw[->] (dec2) -- node[lbl, right] {Yes} (pos);
  \draw[->] (pos) -- (dec3);
  \draw[->] (dec3) -- node[lbl, right] {No} (lyric);
  \draw[->] (lyric) -- (flip);
  \draw[->] (dec3.east) -- node[lbl, above] {Yes} (reset.west);
  \draw[->] (reset.south) |- (lyric.east);
  \draw[->] (dec1.west) -- node[lbl, above] {No} ++(-1.6,0) |- (wait.west);
  \draw[->] (dec2.west) -- node[lbl, above] {No} ++(-1.6,0) |- (wait.west);
  \draw[->] (flip.east) -- (lane |- flip.east) -- (lane |- wait.east) -- (wait.east);
\end{tikzpicture}
\end{center}
\caption{LCD表示子機の処理フロー}
\label{fig:flow_lcd}
\end{figure}
